%\addchap{\lsAcknowledgementTitle} 
\addchap{Preface and acknowledgments}
Early on in this project we made a commitment to ourselves: \it{we will follow the data -- wherever it might lead}.
If it led to supporting certain subgrouping hypotheses prevalent for several decades, so be it.
At l\ili{east} we would then know that there was some substance behind these hypotheses,
and have a clearer idea of what that might be.
If it led in other directions or to other conclusions, that was fine too,
since we would also have a clearer idea of how the Austronesian languages in the Wallacea region came to be what they are today.
We wanted to try to whittle away significantly at the big ``question mark''
that has been hanging over the classification of the Austronesian
languages in our target region of eastern Indonesia and Timor-Leste for 150 years or more.

Early on and even around 90{\%} of the way through the writing of this work,
we did not know what the last chapters would look like.
We had to work through the chapters on individual subgroups first before we could explore ways in which they might be linked together.
In following the data we had several false starts as we tried out different possibilities.
We had to identify and confront our own assumptions,
particularly when they showed themselves to be in conflict with the story the data were telling us.
Our early attempts at maps showing subgroups in the region had to be revised several times along the way -- even fairly late in the drafting.
The organisation of our chapters on different subgroups had to be changed and rearranged several times, even fairly late in the process.
\ili{As} we encountered additional data, we had to go back and revise sections already drafted.
We had several dilemmas over which we conferred: What are the implications of this? If we interpret this as X,
what is the ripple effect? If we interpret this as Y,
what is the ripple effect? Back to our commitment: \it{the data are king; follow the data.
A subgrouping hypothesis must be subservient to the data,
not the other way around.} In at l\ili{east} one case,
we realised we were being inconsistent:
in one place we were treating two different groups as a very weak higher-level
subgroup based on an extremely common and non-exclusive merger,
whereas with two other groups in a nearby region we were saying
such a common and widespread merger was insufficient for subgrouping.
So we went back and revised several chapters (and maps) for consistency.
\ili{As} we encountered new data or drafted new sections,
we started connecting the dots in new and different ways,
and that required modifying sections already drafted in the earlier chapters,
sometimes expanding them to include additional examples,
as well as rethinking how the later chapters should come together.

In seeking to understand the languages in our target region,
we also had to look at languages to the \ili{east}, west,
and north to provide a context for comparison.
In gathering data from many languages and many regions,
appendix \ref{ch:LanLis} lists the sources from which we have gleaned
data from at l\ili{east} 517 languages/speech varieties.
425 are Austronesian and 91 are from various Papuan (non-Austronesian) families.
292 of these varieties are found within our primary
target region of Linguistic Wallacea,
including 36 Papuan Timor-\ili{Alor}-\ili{Pantar} languages,
but excluding \tsc{South Halmahera-West new Guinea},
\tsc{North Halmahera} languages and other
Papuan languages of the Bird's Head and Cenderawasih Bay region.
This volume has 78 figures (including 32 maps),
334 tables comparing different types of data from many languages,
and 195 numbered examples/lists.
These tallies do not count the many in-text examples found throughout the volume.
Appendix \ref{ch:BanTer} provides terms for ‘banana’ from 391 languages/varieties.
Appendix \ref{ch:SurMar} provides marsupial terms from 285 languages/varieties.

This work has been a true collaboration.
Neither of us individually control all the data or literature in all the regions.
Neither of us could have written this book on our own.
Both of us drafted different sections and different chapters.
Both of us provided additional data or fleshed out chapters the other drafted.
Both of us contributed significantly to revisions in all chapters.

Grimes has worked first-hand with languages in the \tsc{\ili{Bima}-Lembata},
\tsc{Timor-Babar}, \tsc{\ili{Central} Timor}, \tsc{Sula-Buru}, and \tsc{Taliabu} subgroups.
He has also worked on \tsc{South Halmahera-West new Guinea} languages of south Halmahera and \ili{Gebe},
as well as Papuan \tsc{Timor-\ili{Alor}-Pantar} and \tsc{North Halmaheran} languages.
Edwards has worked with languages in the \tsc{\ili{Bima}-Lembata},
\tsc{Timor-Babar}, and \tsc{\ili{Central} Timor} subgroups.
Providing the perspective of context from regions just outside of Linguistic Wallacea,
Grimes has worked with the languages of \tsc{South Sulawesi},
and Edwards has worked with the \tsc{Celebic} \tsc{\ili{Bungku}-Tolaki} languages of South\ili{east} Sulawesi.

In some cases, we have been able to confirm and support what other scholars
have said about the classification of languages in the Wallacean region.
In other cases, we have had to conclude that the fuller data to which
we have access are telling a different story than that told by others in the past.
That should not be surprising in a region in which so
little data are available and so little comparative work has been done,
much of it based on the lexicon in a region of great mobility.

Given the severe under-documentation of the 260 or so Austronesian languages
in our target region of eastern Indonesia and Timor-Leste,
we are hugely thankful to so many friends and scholars,
several of whom have generously made available to us spreadsheets,
PDFs, word lists, and computer files of unpublished data that represent months or even years of their hard work.
We are humbled by the trust extended to us.
We are also aware that they are among the many curious people and interested parties
who would really like to have a clearer idea of the prehistory and classification of the Austronesian languages in Wallacea.
(Sources of data for individual languages -- both published and unpublished -- are listed in appendix \ref{ch:LanLis}.
We focus mostly on unpublished works below.)

None of the people listed below are in any way responsible
for errors in the presentation and/or interpretation of their data.
Similarly, it cannot be assumed that they agree with our conclusions.

A massive thank you to Rick Nivens for providing a database of
over 900 Proto-\ili{Aru} reconstructions with supporting reflexes.
\citet{Nivens2017} compiled data from a number of sources including his
own fieldwork and fieldwork by other linguists, mainly: Jock Hughes (many languages),
Benjamin T. Daigle (\ili{Batuley}, \ili{Mariri}), and Antoinette Schapper (\ili{Ujir}).

Special thanks also go to Timothy Usher who made a significant contribution by organising the \ili{Yamdena},
\ili{Fordata} and \ili{Kei} data from the early 20th century sources \citep{Usher2019},
as well as collating the Bomberai survey data \citep{Donohue2010, Narfafan2011} into a single file.
Timothy Usher tabulated the sound correspondences for Bomberai Tip \citep{Usher2018}.
All this provided a head start on understanding how these languages fit into the picture.

In the Maluku region,
David Coward freely shared his unpublished database of \ili{Selaru} lexical data
from that most unusual and fascinating language.
Carrie Beckley shared her unpublished Toolbox lexical database
including example sentences from \ili{Yalahatan} of south central Seram,
a language about which very little has been previously known.
Toni and Heidi Mettler (the latter now deceased) generously
shared their extensive but as yet unpublished \ili{Yamdena} lexical files and grammar sketch.
Mark and Kathy Taber shared their unpublished lexical database.
John Christensen also shared his unpublished lexical database of \ili{Kisar} (\ili{Meher}).
Roland Walker gave Grimes a photocopy of his unpublished \ili{Kowiai} field notes,
lexical data, and grammar sketch. Craig Marshall kindly
provided much needed last-minute data from \ili{Seluwasan} and \ili{Makatian},
including word lists collected in 2012 by Matt Conner,
and supplementary data collected at our request by Leu Maiseka,
all working under Yayasan Pemberdayaan Masyarakat Desa in Tanimbar.

Also in the Maluku region,
several friends and scholars kindly filled in a copy of \citet{GrimesC1990}
by hand and sent a photocopy of their unpublished data to Grimes
for use by him in comparative studies.
These include: Carrie Beckley, \ili{Yalahatan}; Rosemary Bolton, South \ili{Nuaulu};
John and Sylvia Christensen,
\ili{Kisar}; David and Naomi Coward, \ili{Selaru}; Bryan Hinton, \ili{Tugun};
Jock Hughes, \ili{Dobel}; Wyn and Carol Laidig,
\ili{Larike}; Russel Loski, \ili{Geser}-\ili{Gorom}; Craig and Sarah Marshall, \ili{Fordata}; Takashi Matsumura,
\ili{Irarutu}; Frank McCollum, \ili{Liana-Seti}; Toni and Heidi Mettler, \ili{Yamdena}; Rick Nivens,
Southwest \ili{Tarangan}; Kathy Taber, \ili{Luang}; Yushin and Takako Taguchi,
\ili{Alune}; Masahiro Takata, \ili{Kola}; Dirk Teljeur, \ili{Gane}/\ili{Giman}; Ed Travis,
\ili{Kei}; Ron and Jacqui Whisler, \ili{Sawai}.
Ed Travis (now deceased) also provided Grimes with an
\ili{Ambelau} word list from 6 October 1986.
Alistair MacDonald ensured we had a copy of his
Nusawele (\ili{Manusela}) data \citep{MacDonald2015}.

In the greater Timor region,
Thersia Tamelan generously allowed us to cite drafts of her now complete
grammar of \ili{Dela} as well as providing a copy of her lexical database.
Jim Fox has provided \ili{Termanu} data,
as well as a lot of encouragement along the way.
Catharina Williams-van Klinken shared her Toolbox databases
of \ili{Tetun} Dili [tdt] and her earlier \ili{Tetun} [tet] lexical data.
Celicia Elliot generously shared her unpublished \ili{Kemak} word lists,
lexical data and draft grammar sketch.
Native speakers of a number of languages continue to work
with us enthusiastically to compile lexical databases of their languages.
These include: Helena Aplugi and Ayub Ranoh (the latter now deceased),
\ili{Dhao}; Adika Balukh and Albert Zacharias (the latter now deceased),
\ili{Lole}; Roni Bani and Oma Ora,
\ili{Kotos} \ili{Amarasi}; Jusuf Boimau and Adel Timo,
\ili{Amanuban}; Teofigildu Elu and Anita Ulan,
\ili{Baikeno}; Yustin Nako, \ili{Rikou}; João Cristo Rei,
\ili{Galolen}; James Soares, \ili{Raklungu}; Yanri Suana,
Rev. Jumetan (and Kirsten Culhane),
Amfo{\Q}an; Yanti Tunliu-Mooy and Ena Hayer-Pah, \ili{Tii}.

Eline Visser provided us with data on \ili{Uruangnirin} (most available as \citealp{Visser2019b,Visser2023}),
including unpublished grammar notes.
Laura Arnold generously provided us with some unpublished field notes on
marsupial terms from several languages of the Raja Ampat Islands and west New Guinea.
She also provided feedback on a draft of \srf{11-05} (\tsc{Tanimbar-Bomberai}).
Although outside the direct scope of our book,
we also had many stimulating e-mail exchanges with
her on the position and nature of the \tsc{South Halmahera West-New Guinea} languages.

David Mead quite generously provided a huge amount of unpublished data from Sulawesi languages.
Of particular note are many marsupial terms related to several different etyma from dozens of languages.
This has allowed us to see a much fuller picture in the regions of the Austronesian world in which marsupials are found,
and consequently address some of the confusion that has reigned over the reconstruction of marsupial terms in the past.

Anthropologist/sociolinguist Barbara Dix Grimes gave some helpful high-level suggestions
that resulted in adding more maps and re-ordering some sections to communicate more effectively to a broader audience.
Numerous email exchanges with anthropologist James Fox have also helped us clarify how to explain what we were finding to non-linguists.

Malcolm Ross read the first complete draft of the book,
made helpful suggestions and asked insightful questions that
led to searching for more data and ultimately restructuring several chapters once again.
We are grateful for his input.
He is not to blame for the final version.
Similarly, two anonymous referees read the manuscript and provided extensive feedback.
We have carefully considered their comments and have revised the manuscript as appropriate.

We are a product of our histories.
Grimes has benefited from studying Austronesian linguistics with John Wolff at Cornell,
and with Bill Foley at the Australian National University.
With both of us having done our PhDs at ANU,
we have both gained significantly from multiple stimulating interactions with Malcolm Ross and Andy Pawley there.
We have both enjoyed and been challenged by interaction over many years with Mark Donohue.
We have also both enjoyed and benefited from interaction with Marian Klamer at Universiteit Leiden over many years.

We have both also learned a great deal from the prolific writings of and interaction with Bob Blust at University of Hawai‘i.
We regret that the time it took to produce this work meant that it appeared after Blust’s passing,
meaning that we are unable to further engage with him regarding our proposed revision of Austronesian subgrouping.
He was aware of our work,
and we had some e-mail exchanges with him after we presented the initial results of this book at ICAL 15 \citep{EdwardsGrimes2021}.
Although we disagree with his classification of the languages of
eastern Indonesia and Timor-Leste we learnt a lot from him regarding Austronesian historical linguistics,
and our use of the term “Blust line” is our attempt to honour his legacy in this part of the Austronesian world.

In December 1991,
Bernd Nothofer at Frankfurt University posed a question to Grimes,
“If CMP is not a subgroup,
then what is it?” Grimes has been pondering the answer to that question ever since.

In putting together this present study,
we are aware that we have learned to our benefit from scholars who hold views with which we respectfully disagree.
Reviewing their data and addressing their arguments has forced us to sharpen our own thinking,
ferret out additional data,
and try to express more clearly what we think is really going on in the region.
For that we are also grateful.
There is nobody with whom we disagree across the board.
A careful reading of this present work should show that we agree and support other scholars on many issues,
and sometimes disagree with those same scholars on a few other issues.
We follow where the data lead.

We fully expect there to be additional refinements
in the future to the classification presented in this work.
We would not be at all surprised if we ourselves contribute to some of those refinements.
Given the severe under-documentation of languages in the region,
any high-level classification at this fledgeling stage of
our knowledge should be seen as preliminary.
We have tried to lay out the data in ways that others
can see possible patterns that we have not yet seen ourselves,
and we have tried to be transparent about apparent exceptions in the data,
rather than present only the data that are favourable to our conclusions.
It is our hope that this study makes significant inroads into advancing
our collective knowledge and understanding of the prehistory and
classification of the Austronesian languages of eastern Indonesia and Timor-Leste.

Some comments on the timeline regarding this book may also be helpful,
as we are aware that several sources are already citing the work represented in this book.
We first began this project in 2018 when we were invited to contribute a chapter to an edited volume.
We quickly realised that the topic could not be adequately dealt with within the recommended
page limit so we began reorganising the work as a monograph to give the issue its due.
An initial draft was produced in mid-2021,
a summary of which was presented at 15-ICAL \citep{EdwardsGrimes2021}.
Shortly after this, we submitted the manuscript for publication.
Malcolm Ross read the manuscript of the whole book and made extensive
and extremely helpful comments on that first complete draft.
While his comments were mostly positive, they nevertheless
led to significant reorganisation
and scouring for additional data
(including revising the number of subgroups we identified)
which were carried out over the next two years.
The revised manuscript was sent to two anonymous reviewers.
The referee reports were received in February 2024 and revisions
addressing their comments and suggestions were completed in July 2024.
At that point we received a quote far beyond our means for copyediting and typesetting
and were told our intended publisher had a significant backlog.
After several months trying to find a solution,
the original publisher graciously concurred
that we could and should publish elsewhere.
We approached the editorial board of the newly
established \it{Languages of the New Guinea Region}
series at Language Science Press who agreed to publish
the manuscript and provided additional reviewers' comments,
which led to more revision before the eventual final submission.

\begin{flushright}
\it{Charles E. Grimes and Owen Edwards}

\ili{Kupang} and \ili{Ambon}, May 2025
\end{flushright}
